% Generated draft: requires mathematical and source review.
\documentclass[11pt]{article}
\usepackage[margin=25mm]{geometry}
\usepackage{fontspec}
\setmainfont{Latin Modern Roman}
\usepackage{amsmath,amssymb,mathtools,unicode-math}
\setmathfont{Latin Modern Math}
\usepackage{enumitem,needspace,xurl,hyperref}
\setlength{\parindent}{0pt}
\setlength{\parskip}{5pt}
\setlength{\emergencystretch}{3em}
\newcommand{\N}{\mathbb{N}}
\newcommand{\Z}{\mathbb{Z}}
\newcommand{\Q}{\mathbb{Q}}
\newcommand{\R}{\mathbb{R}}
\newcommand{\C}{\mathbb{C}}
\newcommand{\F}{\mathbb{F}}
\newcommand{\A}{\mathbb{A}}
\newcommand{\PP}{\mathbb{P}}
\newcommand{\E}{\mathbb{E}}
\newcommand{\eps}{\varepsilon}
\newcommand{\dd}{\,\mathrm{d}}
\begin{document}
\section*{UnsolvedMath ID: 30006406}
\textbf{Source problem code:} catalogue entry
\subsection{Definitions and mathematical statement}
Let \((M_t)_{t\geq0}\) denote the level-set mean-curvature flow in the supplied record. The ambient space, admissible initial data, and construction of the associated inner and outer flows are not specified. Consequently, the following records the proposed equivalence with its unresolved qualifications; it is not a fully specified conjecture.

In the hypersurface interpretation, let \(X\) denote the ambient space and let \(n\) denote the hypersurface dimension. These roles for \(X\) and \(n\) require confirmation from the source. Nonfattening means that the evolving hypersurface level set has empty ambient interior:
\[
\operatorname{Int}_X(M_t)=\varnothing\qquad\text{for every }t\geq0.
\]
This definition treats \(M_t\) as the hypersurface level set, rather than an evolving enclosed region. Write \(M_t^{\mathrm{in}}\) and \(M_t^{\mathrm{out}}\) for the associated inner and outer hypersurface flows, once their construction and initial-region conventions have been fixed. Nondiscrepancy is the condition
\[
M_t^{\mathrm{in}}=M_t=M_t^{\mathrm{out}}
\qquad\text{for every }t\geq0.
\]
The supplied record does not provide enough information to define these two associated flows independently.

The stated intersection principle quantifies over every smooth closed mean-curvature flow \((N_t)_{t\geq0}\) in the same ambient space satisfying
\[
N_t\not\subseteq M_t\qquad\text{for every }t\geq0.
\]
Here a smooth closed flow is understood to consist of compact hypersurfaces without boundary evolving with normal velocity equal to their mean-curvature vector. The displayed noncontainment is an explicitly provisional reading of the corrupted source expression \(N_t\,6\subseteq M_t\). For each such comparison flow, the principle requires
\[
\dim(M_t\cap N_t)\leq\dim(M_s\cap N_s)
\qquad(0\leq s\leq t),
\]
and also
\[
\forall s\geq0,\qquad
\dim(M_s\cap N_s)<n-1
\ \Longrightarrow\
\bigl(\forall t>s,\ M_t\cap N_t=\varnothing\bigr).
\]
The meaning of \(\dim\), its value on the empty set, and the meaning of \(n\) are not specified in the record. These dimension statements therefore require those conventions before they become precise. The all-time comparison-flow quantifier is retained as supplied; replacing it by a common smooth existence interval would require source confirmation.

The problem is to determine under which hypotheses, and in which directions, the following equivalences hold:
\[
\text{nonfattening}
\ \stackrel{?}{\Longleftrightarrow}\
\text{nondiscrepancy}
\ \stackrel{?}{\Longleftrightarrow}\
\text{the stated intersection principle}.
\]
The original excerpt attributes a special-case result to Hershkovits--White for flows with ``mean convex neighborhoods of'', but the sentence ends at that point. The missing hypothesis and the exact scope of that attribution cannot be recovered from the supplied text.
\subsection{Short English statement}
Determine when nonfattening, agreement of the inner and outer flows, and the stated intersection principle characterize the same level-set mean-curvature flows. The supplied record labels the general question open, but missing geometric hypotheses, dimension conventions, and a truncated special-case attribution prevent a definitive formulation without source review.
\subsection{Sources}
\begin{itemize}
\item Supplied local UnsolvedMath registry record, permanent ID \(30006406\), title: Nonfattening and Intersection Principles for Level-Set Flows; exact source problem code: \texttt{catalogue entry}. No individual proposer is supplied; the proposed year is recorded as 2025.
\item Supplied report citation: \emph{Partial Differential Equations (2026)}, Oberwolfach Reports, DOI 10.4171/owr/2025/33, pp.\ 1753--1822. \url{https://doi.org/10.4171/owr/2025/33}. The record describes Oberwolfach Reports as documenting mathematical workshops and their open-problem discussions. The citation's title year 2026 differs from its DOI year and the recorded proposed year 2025; this discrepancy is unresolved.
\item \emph{An Intersection Principle for Mean Curvature Flow}. Supplied literature reference: \url{https://arxiv.org/abs/2505.11600}.
\item \emph{Non-fattening of Mean Curvature Flow at Singularities of Mean Convex Type}. Supplied literature reference: \url{https://arxiv.org/abs/1704.00431}. The original excerpt separately attributes a special-case result to Hershkovits--White, with an incomplete hypothesis.
\item Literature assessment dated exactly \texttt{2026-08-22}: the supplied curation labels the general equivalence between nonfattening and the intersection principle unresolved and describes the 2025 intersection-principle work as proving important finite-singularity and conditional cases. This is supplied dated curation derived from the OWR record, not an independently verified literature assessment.
\item The registry labels the statement \texttt{corrected\_verified} and says that the official report was checked, the question isolated from surrounding exposition, and its notation repaired. These are supplied verification claims; no source visit or independent mathematical verification is asserted here.
\end{itemize}
\end{document}
